% Transcription — fonds Alexandre Grothendieck, shelfmark 76, batch 3 (pages 41-60).
% Established from the facsimile archives/batches/76/batch-03.pdf, cut from a
% copy of 76.pdf fetched through the project's relay
% (https://grothendieck-relay.onrender.com/source/76.pdf); PDF page k =
% archivists' page 40+k, no cover sheet. Per the skill transcribe-grothendieck.
% The folder is sixty-two pages in four batches (1-20, 21-40, 41-60, 61-62),
% transcribed in parallel; batches 2 and 4 were read for notation and
% alignment once they existed (p. 41 continues batch 2's run on the
% categories F_n and the functor psi_n; p. 60's « Prop » is the one batch 4's
% two corollaries on p. 61 follow).
% Pass: Opus 5.5 (claude-opus-5-5), 2026-09-26 — first pass, unchecked against the pages by a human.
%
% Dating and title. The dating is the folder's own, copied verbatim from the
% catalogue; the group « Géométrie et topologie combinatoire » (69-89) holds
% it, and since the folder has a date of its own the group's range is not
% added. The title is the catalogue's, without its final full stop.
%
% Pages skipped: 44 and 48 (blank), 46 (blank but for show-through of p. 47).
% Page 43 is a cover in his hand (« 2-cartes cellulaires » struck, « cube
% combinatoire » above; « (16 p) » in pencil, perhaps not his): no \page{},
% its words are the section heading, with a note. Pages summarised: none;
% nothing is typed and nothing is by another hand.
%
% Pages transcribed: 41-42 (blue ink on perforated listing paper, the end of
% the programme of numbered theorems begun in batch 2), 45 (listing paper:
% small sets as finite geometries over F_2, F_3, F_4, F_5; card G(F_q) for
% Chevalley groups via the cells of G/B, with a hand-drawn table), 47
% (« Sujets DEA », four numbered subjects), 49-60 (black felt-tip, the cube
% run). The hand is fast but mostly legible in the cube run; the connective
% prose of pp. 45, 55-56 and 58 is the weakest.
%
% His pagination: none on the leaves. Numbered runs: Th. 5, 6, 7 (pp. 41-42)
% continue his Th. 1-4 of batch 2; « Sujets DEA » 1)-4) (p. 47); on p. 57 a
% struck « 1°) »; p. 58-59 cases a) and b) (« cas tordu »); p. 59 example
% with 1°)-2°) and circled (+), (-); p. 60 two runs 1°)-4°) (under b), two
% cycles, and c), three cycles), each case carrying in the left margin a
% sign and a count (« -, 6 », « +, 1 » ...), recorded as \marginal. No
% « cf. p. N » cross-reference; « cf cor. précédent » on p. 50.
%
% What the batch is about. Pp. 41-42 close the programme on cellular n-maps:
% barycentric subdivision of psi_n(C_n) (Th. 5), the simplicial complex of
% the twice-subdivided triangulation in terms of the cascade (Th. 6), the
% essential image of F_n° in combinatorial terms, linked to recognising
% spheres and « l'hyp. de Poincaré » (Th. 7), and further points (flags,
% duality, products, fibrations). P. 45 is an aside. From p. 49 the run
% « cube combinatoire »: the cube C(B, sigma) of a finite set B with a
% fixed-point-free involution, E = Hom_{±1}(B, R); faces <-> partial
% sections of B over I = B/sigma (order-reversing); isomorphisms of
% (B, sigma) = affine isomorphisms of cubes = isomorphisms of their face
% posets; Aut(C) simply transitive on flags (the lemma: an automorphism of a
% convex polytope fixing a flag is the identity), cubes <-> right
% G_0-torsors, G_0 = {±1}^n · S_n; the antipody and the centre of Aut(C);
% f-vector c_d = binom(n,d) 2^{n-d} and Euler-Poincaré; orientations
% (sigma acts by (-1)^n; a combinatorial definition of orientation via
% « repères »); regular and homogeneous polytopes (N | N_c | R); conjugacy
% classes of Aut(C_n) by two sequences alpha_i^±, the fixed faces
% (generating function prod (1+2t^i)^{alpha_i^+}), the ten classes for
% n = 3; and the « Prop » (u x = ±x for all vertices => u = ±id, n >= 3)
% whose corollaries open batch 4.
%
% Notation fixed for the batch, following batches 1, 2 and 4: \mathbb{B} for
% his double-stroked B; \Phi for his barred I (Phi_0(C), Phi_{≤x}); Fac(C)
% and Comb(C) for the face poset as he writes it; Rep(C), Drap(C);
% \mathfrak{S} for the symmetric group; \mathcal{F}_n, (Casc)_n, psi_n,
% \widetilde{X}_i as in batch 2.
%
% Apparatus: 81 \ill{}, 64 \uncertain{}. Hardest pages 45 (dense, two unrelated blocks),
% 55-56 (the connectivity argument and the counting), 58 (the tableau and
% the right-hand note).

\documentclass[11pt,a4paper]{article}
% Le préambule commun à toutes les transcriptions du fonds Grothendieck.
%
% Il tient en une idée : **l'incertitude se compose, elle ne se résout pas.**
% Une transcription qui lisse un mot douteux a détruit la seule chose pour
% laquelle on la faisait — un lecteur doit pouvoir savoir, à chaque endroit,
% ce qui était sur le feuillet et ce qui a été ajouté. Les six macros qui
% suivent sont donc l'appareil critique, et non de la décoration.
%
% Les mêmes commandes sont reconnues par `scripts/render.mjs`, qui produit la
% vue de lecture du site. Une macro ajoutée ici doit l'être aussi là-bas,
% faute de quoi le rendu échouera bruyamment — ce qui est voulu : un rendu
% silencieusement faux est pire qu'un rendu absent.

\ProvidesPackage{grothendieck}[2026/08/08 Transcriptions du fonds Grothendieck]

\usepackage{amsmath,amssymb,amsthm}
\usepackage{mathrsfs}
\usepackage{stmaryrd}
% Les diagrammes commutatifs sont partout dans ce fonds : c'est la langue
% même de la géométrie algébrique de Grothendieck, et une transcription qui
% les rendrait en prose ne transcrirait rien.
\usepackage{tikz-cd}
% Français par défaut — c'est la langue des feuillets — mais l'anglais est
% chargé aussi : la traduction et le résumé sont des documents anglais, et la
% typographie française (espaces avant les deux-points) y serait une faute.
% Ces documents basculent par \selectlanguage{english} en tête de corps.
\usepackage[english,french]{babel}
\usepackage{fontspec}
\usepackage{geometry}
\geometry{a4paper,margin=25mm,marginparwidth=28mm}
\usepackage{marginnote}
\usepackage{xcolor}
% ulem, not soul. `soul`'s \st and \ul are built on a character-by-character
% scan that XeLaTeX's font machinery breaks: the document fails with an
% undefined control sequence rather than degrading. ulem does the same two jobs
% and is Unicode-engine safe. `normalem` stops it hijacking \emph, which the
% transcriptions use for Grothendieck's own emphasis.
\usepackage[normalem]{ulem}
\usepackage{hyperref}
\hypersetup{colorlinks=true,linkcolor=black,urlcolor=black,pdfborder={0 0 0}}

% --- Métadonnées du lot -----------------------------------------------------
% Reprises telles quelles par le script de rendu, qui les affiche en tête de
% la vue de lecture. Elles fixent ce que le fichier prétend couvrir : sans
% elles, une transcription flotte sans point d'attache dans un fonds de seize
% mille feuillets.

\newcommand{\folder}[1]{\gdef\grothFolder{#1}}
\newcommand{\batch}[1]{\gdef\grothBatch{#1}}
\newcommand{\leaves}[2]{\gdef\grothFirst{#1}\gdef\grothLast{#2}}
\newcommand{\foldertitle}[1]{\gdef\grothFoldertitle{#1}}
\gdef\grothFolder{?}\gdef\grothBatch{?}\gdef\grothFirst{?}\gdef\grothLast{?}\gdef\grothFoldertitle{}

% --- Le résumé --------------------------------------------------------------
%
% Il ouvre la lecture modernisée, sous le titre « Résumé », et il est composé
% en petit. Ce n'est pas un ornement : c'est le seul passage du document qui
% ne soit pas des mathématiques, et un lecteur doit pouvoir le reconnaître
% avant de l'avoir lu — puis le sauter s'il connaît déjà le sujet, ou s'y
% arrêter s'il ne le connaît pas. Le filet qui le suit marque la reprise du
% corps.

\newenvironment{resume}
  {\small\setlength{\parskip}{0.45em}}
  {\par\medskip\hrule\medskip}

% --- Mots-clés ---------------------------------------------------------------
%
% En anglais, en clôture du résumé de la lecture modernisée : le vocabulaire
% moderne sous lequel chercher ce que les feuillets construisent. Le manifeste
% du site (`npm run manifest`) les extrait pour étiqueter la cote entière —
% cette ligne est la source unique des tags, il n'y a pas de fichier à côté.
\newcommand{\keywords}[1]{\par\smallskip\noindent
  {\footnotesize\textsc{Keywords} --- \textit{#1}}\par}

% --- L'appareil critique ----------------------------------------------------

% Le numéro de feuillet, en marge. C'est l'ancre qui permet de régler un
% désaccord sur une formule en regardant *un* feuillet plutôt qu'une chemise
% de six cents. Le site s'en sert aussi pour tourner les pages du fac-similé
% quand on fait défiler la transcription.
\newcommand{\leaf}[1]{%
  \marginnote{\footnotesize\color{gray!70}\textsf{#1}}%
  \label{leaf:#1}\ignorespaces}

% Illisible : on ne devine pas. Un mot inventé qui se lit comme les autres est
% le pire résultat possible.
\newcommand{\ill}{\textcolor{red!60!black}{[\,\ldots\,]}}

% Lecture douteuse : le mot est donné, et signalé comme incertain.
\newcommand{\uncertain}[1]{\uline{#1}}

% Ajout de l'éditeur — une lettre manquante, un mot sous-entendu.
\newcommand{\add}[1]{\textcolor{blue!50!black}{[#1]}}

% Biffé par Grothendieck lui-même. Ce qu'il a rayé fait partie du document :
% une rature dit souvent où la pensée a changé de direction.
\newcommand{\struck}[1]{\sout{#1}}

% Note du transcripteur, sur l'état matériel du feuillet.
\newcommand{\note}[1]{\par\smallskip{\small\color{gray!60!black}\textit{[#1]}}\par\smallskip}

% Note marginale de Grothendieck, distincte du corps du texte.
\newcommand{\marginal}[1]{\par\smallskip\noindent\hspace{1em}%
  {\small\color{gray!60!black}#1}\par\smallskip}

% --- Titre ------------------------------------------------------------------

\AtBeginDocument{%
  \begin{center}
    {\large\bfseries Fonds Alexandre Grothendieck --- cote n\textsuperscript{o}~\grothFolder}\\[2pt]
    {\itshape\grothFoldertitle}\\[4pt]
    {\small feuillets \grothFirst--\grothLast\ (lot \grothBatch)}\\[2pt]
    {\footnotesize\color{gray!60!black}Transcription établie d'après le fac-similé numérisé
     par l'Université de Montpellier.\\
     Les lectures incertaines sont soulignées, les illisibles notées [\,\ldots\,],
     les ajouts entre crochets.}
  \end{center}
  \vspace{1em}}

\endinput


\folder{76}
\batch{3}
\pages{41}{60}
\dating{[vers 1977]}
\watermark{Édition de démonstration}
\foldertitle{n-cartes cellulaires : notes manuscrites (s.d.)}

\begin{document}

\page{41}
\note{la page continue le programme de théorèmes numérotés de la page 40
(Th.\ 1 à 4, construction de $\psi_n : (\mathrm{Casc})_n \to
\mathcal{F}_n$ par « multicône ») ; les numéros 5, 6, 7 sont les siens}

Enfin, on devrait donner un

\emph{Th.\ 5} Subdivisions barycentriques de $\psi_n(C_n)$, et description
par ens.\ \uncertain{pseudo-simpliciaux}\ldots
\note{un trait vertical dans la marge gauche accompagne ces deux lignes}

\emph{NB} La donnée d'un homéom.\ $X_n \simeq \psi_n\varphi_n(X_n)$
revient à la donnée, pour tout $i$ et toute composante connexe
$\widetilde{X}_i^\alpha$ de $\widetilde{X}_i$ (« $i$-cellules ») d'un
homéomorphisme
\[
  \widetilde{X}_i^\alpha \xrightarrow{\ \sim\ }
  \mathrm{C\hat{o}n}(\partial\widetilde{X}_i^\alpha)
\]
induisant l'identité sur le bord. [La subdivision barycentrique se fait
\uncertain{par dessus} une triangulation. On pourrait expliciter la
subdivision barycentrique itérée 2 fois, qui correspond, elle, à une
triangulation. Il faudrait donc expliciter le \uncertain{complexe}
simplicial (ensemble $+$ ens.\ de parties finies, stable par inclusion) de
cet espace triangulé, en termes de la cascade $C_n$, et donner une
description \ill{} \ill{} \ill{} de la filtration \ill{} \ill{} sa
réalisation topologique. Cela fera bien encore un th.\ 6 (et peut-être un
th.\ 7\ldots)
\note{« par dessus » : lecture incertaine ; le crochet ouvrant est tracé
en équerre, sans crochet fermant sur la page}

\page{42}
Je veux dire : caractérisons l'image essentielle de
$\mathcal{F}_n^\circ$ (\uncertain{$C_n \to$ ens.}) [en termes
combinatoires ! Voire, ce qui revient au même : étudier si un espace
triangulé est homéom.\ à une sphère --- \struck{mais déjà que \ill{}}
\add{cela semble lié : l'hyp.} \add{de Poincaré !}] ; on en conclura aussi
la caractérisation de l'image essentielle de $\mathcal{F}_n \to C_n$. On
aura ainsi un

\emph{Th.\ 7}
\note{ici « Th.\ 7 » est doublement souligné et n'est suivi d'aucun
énoncé ; la « lettre » $\mathcal{F}_n$ est son F gras, avec un petit
cercle en exposant la première et la dernière fois ; « cela semble lié :
l'hyp.\ de Poincaré ! » est écrit au-dessus de la ligne, sur les mots
biffés}

Autres points : examiner
\begin{itemize}
  \item[$*$] Relations avec ensembles ordonnés (polyèdres combinatoires)
  via drapeaux\ldots
  \item[$*$] Dualité dans $\mathcal{F}_n^\circ$
  \item[$*$] Produits d'espaces : filtrations admissibles
  \item[$*$] Fibrations ?
\end{itemize}

\section*{Cube combinatoire}
\note{titre tiré de la page 43, une feuille de couverture de sa main :
« \struck{2-cartes cellulaires} » biffé, remplacé par « cube combinatoire
» ; en haut à gauche de la même feuille, au crayon et entre parenthèses,
« (16 p) », qui pourrait ne pas être de sa main}

\page{45}
\note{feuille de listing perforée, encre bleue ; la page porte deux
blocs sans lien apparent avec les cubes : en haut des petits ensembles
identifiés à des géométries finies, en bas le dénombrement de groupes
linéaires sur $\mathbf{F}_q$}

\emph{Ens de card 2} : \uncertain{torseur} sous $\mathbf{F}_2$ i.e.\
\add{droite} \struck{\ill{}} affine \struck{\ill{}} $\mathbf{F}_2$

\emph{Ens de card 3} : droite affine sur $\mathbf{F}_3$ / éléments non nuls
d'un vectoriel de dim 2 sur $\mathbf{F}_2$ / droite projective sur
$\mathbf{F}_2$ / $+$ orientation : \struck{\ill{}}

\emph{Ens de card 4} : \uncertain{plan} affine sur $\mathbf{F}_2$ / droite
projective sur $\mathbf{F}_3$ / \emph{orientation} : droite affine
\uncertain{sur} $\mathbf{F}_4$ ?

\emph{Ens de card 5} ($\leftrightarrow$ droite projective sur
$\mathbf{F}_5$, les 6 pts \ill{} les \ill{} des \ill{}
\uncertain{icosaèdre} \ill{} à $E$)

\emph{Ens orientés de card 5} droites projectives sur $\mathbf{F}_4$ ---
l'ens.\ :
\note{la ligne s'arrête sur les deux points ; sur la ligne « card 3 », un
court mot est souligné sous « orientation » ; la fin de la ligne « card 5
» est très serrée}

Cas $G = GP(2, \mathbf{F}_q)$, ($N = 3$, \add{$r = 2$} $d_\lambda = 2, 3$)
\[
  \boxed{\ \mathrm{card}\, G(\mathbf{F}_q) = q^N (q-1)^r
  \prod_{1 \leqslant \lambda \leqslant r}
  \frac{q^{d_\lambda} - 1}{q - 1}\ }
\]
\note{au-dessus du produit, une accolade renvoie à « \uncertain{légende}
$q^{d_\lambda - 1}$ » ; « $r = 2$ » est ajouté au-dessus de la parenthèse}
\[
  P_{G/B} = \frac{(1 - T^2)^r}{\prod_\lambda (1 - T^{2d_\lambda})}
\]
\[
  Sl(3, \mathbf{F}_q) \quad \mathrm{card} \ \sum_i b\, q^{d(i)} \qquad
  b = q^3 (q-1)^2 \sum_{0 \leqslant i \leqslant 3} b_{2i}\, q^i
\]
\[
  \Bigl(1 : P_{G/B}(\sqrt{q})\Bigr) = \frac{(1-q)^r}{\prod_\lambda (1 -
  q^{d_\lambda})} = \frac{(q-1)^r}{\prod_\lambda (q^{d_\lambda} - 1)}
\]
\[
  \mathrm{card}\, G(\mathbf{F}_q) = q^N (q-1)^r\, \mathrm{card}\,
  G/B(\mathbf{F}_q)
\]
\note{sous « $G(\mathbf{F}_q)$ », une ligne reliée par un trait :
« \uncertain{gr.\ de Chevalley} » ; au-dessous : « $N$ = nb de racines
\uncertain{pos.}\ $= \dim G/B = \dim B$\,\ill{} », « $r = \mathrm{rg}\,G$ » ;
un premier facteur devant $q^N$ est biffé}
\[
  \mathrm{card}\, G/B(\mathbf{F}_q) = \sum_{\substack{\alpha \text{
  cellules} \\ (d(\alpha) = \text{dim})}} q^{d(\alpha)} = \sum_{0
  \leqslant i \leqslant N} b_{2i}\, q^i = P_{G/B}(\sqrt{q})
\]
\[
  B \supset T, \quad G/T \simeq G/B \qquad P_{B_G} \cdot P_{G/B} = P_{B_T}
\]
\note{sous la dernière égalité : $\bigl(\frac{1}{1-T^2}\bigr)^r$ sous
$P_{B_T}$, relié par une double flèche à $\prod_\lambda
\frac{1}{1-T^{2d_\lambda}}$ sous $P_{B_G}$ ; un mot court sur la flèche,
illisible ; les indices $B_G$, $B_T$ (espaces classifiants) sont des
lectures \uncertain{incertaines}}

\[
  \begin{array}{c|cccc|l}
    q & 2 & 3 & 4 & 5 & \\
    \hline
    & 1 & 2 & 3 & 4 & Gl(1, \mathbf{F}_q) \\
    & 2 & 6 & 12 & 20 & \mathrm{Aff}(1, \mathbf{F}_q) \\
    & 6 & 24 & 60 & 120 & Sl(2, \mathbf{F}_q) \\
    & 6 & 24 & 60 & 120 & GP(1, \mathbf{F}_q) \\
    & 6 & 48 & 180 & 480 & Gl(2, \mathbf{F}_q) \\
    & 24 & 9\cdot 48 & 16 \cdot 180 & 25 \cdot 480 & \mathrm{Aff}(2,
    \mathbf{F}_q)
  \end{array}
\]
\note{le tableau est tracé à la main ; la dernière ligne est écrite à
part, à droite du tableau, et « $16 \cdot 180$ » est souligné ; dans la
première ligne, le nom du groupe est surchargé (un premier « $Sl$ »
récrit) ; en marge des lignes $Sl(2)$ et $GP(1)$, une flèche marquée
« \uncertain{1, 2} »}
\begin{align*}
  Gl(1, \mathbf{F}_q) &\simeq \mathbf{F}_q^* \simeq \mathbf{Z}/(q-1),
  && \mathrm{card}\ q - 1 \\
  \mathrm{Aff}(1, \mathbf{F}_q) &\simeq \mathbf{F}_q^* \cdot \mathbf{F}_q
  && \mathrm{card}\ q(q-1) \\
  GP(1, \mathbf{F}_q) &\simeq Gl(2, \mathbf{F}_q)/\mathbf{F}_q^*
\end{align*}
\[
  \mathrm{card}\ Sl(2, \mathbf{F}_q) = \sum_{\text{cell.\ } i} b\,
  q^{d(i)} \qquad b = q(q-1) \qquad \mathrm{card}\ Sl(2, \mathbf{F}_q) =
  q(q-1)(q+1)
\]
\[
  \mathrm{card}\ GP(1, \mathbf{F}_q) = (q+1)q(q-1), \qquad \mathrm{card}\
  \mathrm{Aff}(2, \mathbf{F}_q) = q^3 (q-1)^{2} (q+1)
\]
\[
  \mathrm{card}\ Gl(2, \mathbf{F}_q) = q(q-1)^2(q+1)
\]
\note{dans « card $Sl(2, \mathbf{F}_q) = \sum$ », « card » est ajouté
au-dessus ; dans « card $GP(1)$ » le premier facteur est surchargé
($q+1$ récrit) ; l'exposant $2$ de $(q-1)$ dans card $\mathrm{Aff}(2)$
est une lecture \uncertain{incertaine} (un petit trait)}

\page{47}
\emph{Sujets DEA}
\begin{enumerate}
  \item[1)] Catégorie isotopique des 1-compl.\ top.\ cpcts avec groupe
  fini d'automorphismes (et plus de morphismes que les seuls
  isomorphismes)
  \item[2)] Catégorie isotopique des cartes cellulaires
  \item[3)] Examiner la notion de convexité du point de vue de la
  géométrie projective. Question : pour $A$ convexe dans un espace
  projectif réel, y a-t-il un hyperplan qui ne le rencontre pas ?
  \item[4)] Géométrie des cubes et espaces homogènes de
  $\mathfrak{S}_4$.
\end{enumerate}

\page{49}
$(\mathbb{B}, \sigma)$ \quad $\mathbb{B}$ ens.\ fini, \struck{\ill{}}
involution dans $\mathbb{B}$ sans pt fixe
\[
  \downarrow
\]
\begin{align*}
  E &= E(\mathbb{B}, \sigma) = \mathrm{Hom}_{\lbrace \pm 1
  \rbrace}(\mathbb{B}, \mathbf{R}) \\
  C &= C(\mathbb{B}, \sigma) = \lbrace \varphi \in E \mid |\varphi(j)|
  \leqslant 1 \ \forall j \in \mathbb{B} \rbrace
  \qquad \mathbb{B} \hookrightarrow C \hookrightarrow E
\end{align*}

\emph{Proposition} Pour toute section partielle $I' \subset \mathbb{B}$ de
$\mathbb{B}$ sur $I = \mathbb{B}/\sigma$, soit
\[
  F(I') = \lbrace \varphi \in C \mid \varphi(j) = 1 \ \forall j \in I',\
  \varphi(j) \neq \pm 1 \ \forall j \in \mathbb{B} - I' - \sigma I'
  \rbrace
\]
Alors l'application $I' \mapsto F(I')$ est une application bijective,
renversant l'ordre
\[
  \mathrm{Spd}(\mathbb{B}/I) \xrightarrow{\ \sim\ } \bigl(\mathrm{Fac}(C)
  - \lbrace \emptyset \rbrace\bigr)
\]
\note{la seconde condition de $F(I')$ est écrite sous la première, après
un mot biffé ; « Spd » (sections partielles) est tel qu'il l'écrit ; sous
« Fac », un mot biffé illisible, et « Fac » lui-même est surchargé}

Ceci précise de façon simple le polyèdre combinatoire associé au polyèdre
$C = C(\mathbb{B}, \sigma)$ en termes de $(\mathbb{B}, \sigma)$, comme
étant (\uncertain{moyennant} l'addition d'un plus petit élément) l'ens.\
ordonné opposé de l'ens.\ des parties $I'$ de $\mathbb{B}$ telles que $I'
\cap \sigma I' = \emptyset$.

\marginal{\emph{NB} Dans le cas du cube ($p$ de degré 2 ?) les él.\ de
$\mathbb{B}$ correspondent aux facettes de codim 1, et pour une telle
facette $F$, $\sigma F$ est l'unique facette de codim 1 qui ne rencontre
pas $F$ (en fait, les hyperpl.\ d'appui correspondants sont
\uncertain{parallèles})}
\note{note écrite verticalement dans la marge gauche, séparée du texte par
un trait ; « de codim 1 » est ajouté au-dessus de la ligne}

\emph{Proposition} Soit $\mathbb{B} \xrightarrow{p} I$ une application
surjective, et soit $\mathcal{P}(p)$ l'ens.\ ordonné par inclusion des
sections partielles de $\mathbb{B}$ sur $I$. Alors $p \mapsto
\mathcal{P}(p)$ est un foncteur pl.\ fidèle du groupoïde des surjections
\struck{\ill{}} (on a comme morphismes les isom.
\[
  \begin{array}{ccc}
    \mathbb{B} & \xrightarrow{\ \sim\ } & \mathbb{B}' \\
    \downarrow p & & \downarrow p' \\
    I & \xrightarrow{\ \sim\ } & I'
  \end{array}
\]
) dans la catégorie des ens.\ ordonnés. On reconstitue
\[
  \mathbb{B} \hookrightarrow \mathcal{P}
\]
ens.\ des pts \struck{\ill{}} minimaux de $\mathcal{P} - \lbrace \emptyset
\rbrace$, et la relation d'équivalence définie par $p$ par
\[
  p(b) = p(b') \Longleftrightarrow (b = b' \text{ ou } C_b \cap C_{b'} =
  \emptyset),
\]
où $\forall b \in \mathbb{B}$, $C_b$ est \struck{l'ens.\ des éléments}
\add{le sous-ens.\ de l'ens.\ $S$} maximaux de $\mathcal{P}$ ($S \simeq
\prod_{i \in I} \mathbb{B}_i$) \struck{\ill{}} qui majorent $b$.
\note{il écrit $\mathcal{P}(p)$ avec un P d'écriture ; le mot biffé devant
« minimaux » est illisible ; la correction « sous-ens.\ de l'ens.\ $S$ »
est écrite au-dessus de « des éléments », qu'elle remplace}

\page{50}
\struck{Ainsi, si $S$ est l'ens.\ des sommets du cube $C(\mathbb{B},
\sigma)$ --- correspondant aux sections de $\mathbb{B}$ sur $I =
\mathbb{B}/\sigma$, \ill{} \ill{} $\in S$ \ill{} l'intersection des
facettes}
\note{ce paragraphe est encadré et barré de traits obliques}

\emph{Corollaire} Soient $(\mathbb{B}, \sigma)$, $(\mathbb{B}', \sigma')$
deux \struck{\ill{}} \add{\ill{}} \ill{}, \struck{Alors} définissant les
cubes $C$, $C'$, on sait que l'on a
\[
  \mathrm{Isom}\bigl((\mathbb{B}, \sigma), (\mathbb{B}', \sigma')\bigr)
  \xrightarrow{\ \alpha\ } \mathrm{Isom}_{\mathrm{aff}}(C, C')
  \xrightarrow{\ \beta\ } \mathrm{Isom}_{\mathrm{ord}}\bigl(\mathrm{Comb}(C),
  \mathrm{Comb}(C')\bigr)
\]
Ces applications sont des bijections. En particulier, on a
\[
  \mathrm{Aut}\bigl((\mathbb{B}, \sigma)\bigr) \xrightarrow{\ \sim\ }
  \mathrm{Aut}_{\mathrm{aff}}(C) \xrightarrow{\ \sim\ }
  \mathrm{Aut}_{\mathrm{ord}}\bigl(\mathrm{Comb}(C)\bigr)
\]
\[
  \begin{array}{l}
    \quad \| \\
    \text{commutant de } \sigma \text{ dans } \mathfrak{S}_{\mathbb{B}} \\
    \quad \cdots \\
    \mathfrak{S}_I \cdot \lbrace \pm 1 \rbrace^I
  \end{array}
\]
\note{un trait vertical dans la marge gauche accompagne l'énoncé ; entre
« commutant » et $\mathfrak{S}_I \cdot \lbrace \pm 1 \rbrace^I$, une ligne
courte (« \ill{} com.\ »)}

\emph{Dém} On sait déjà que $\beta\alpha$ est un isomorphisme (cf cor.\
précédent). Je veux : prouver que $\beta$ l'est. On en sait déjà que
$\beta$ est surjectif ($\beta\alpha$ bij), or il est injectif (car un isom
$C \to C'$ est connu quand on connaît son effet sur les seuls pts
extrémaux) --- cqfd.

\emph{NB} Au lieu de $\mathrm{Isom}_{\mathrm{aff}}(C, C')$ on peut aussi
écrire $\mathrm{Isom}_{\mathrm{proj}}(C, C')$

\marginal{\ill{}}
\emph{Corollaire} Les facettes d'un cube sont des cubes
\note{dans la marge, devant une accolade, un mot court illisible ; la
phrase s'arrête au bas de la page}

\page{51}
\emph{Prop} Le groupe $\mathrm{Aut}(C)$ est simplement transitif sur
$\mathrm{Drap}(C)$ (qui est donc un torseur sous $\mathrm{Aut}(C)$)
\note{« Prop » est surchargé sur un premier mot ; un trait vertical dans
la marge accompagne l'énoncé}

(On dit que $C$ est un pol.\ (convexe) \emph{régulier} pour dire qu'il en
est ainsi\ldots)

\struck{$\mathrm{Card}(\mathrm{Drap}(C)) = 2^n \cdot 2^{\cdots}$}
\note{l'exposant du second facteur biffé est illisible : \ill{}}

Le nb de facettes de dim $d+1$ \uncertain{contenant} une facette donnée
de dim $d$ est \struck{\ill{}}$(n-d)$, donc le nb des \add{drap.}\
\struck{drapeaux} est
\[
  2^n \cdot n \cdot (n-1) \cdots 2 \cdot 1 = 2^n n!
\]
qui est aussi l'ordre de $\mathrm{Aut}(C)$. Il suffit de prouver que
$\mathrm{Aut}(C)$ opère librement sur $\mathrm{Drap}(C)$, ce qui est vrai
pour tt polyèdre convexe :
\note{le signe entre $d$ et $1$ est peu lisible, lu $+$ d'après le compte
$(n-d)$}

\emph{Lemme} Soit $C$ un polyèdre convexe, $u$ un autom.\ de $C$ fixant un
drapeau de $C$. Alors $u$ est l'identité.

Récurrence sur dim $C$. \uncertain{Pour} $\dim C = 0$, c'est
\uncertain{trivial}. Soit $D = (F_0, \dots, F_{n-1}, F_n = C)$ le drapeau
fixé. Par hyp.\ de récurrence, $u$ est l'identité sur $F_{n-1}$, car il
fixe un \add{\uncertain{drap.}} \struck{drapeau} $(F_0, \dots, F_{n-2},
F_{n-1})$. Donc $u$ est un automorphisme affine de $E_C = E$ qui
\struck{\ill{}} \add{induit} l'identité sur un hyperplan $H = E_{F_{n-1}}$,
et qui est d'ordre fini --- or ces seules familles \ill{} c'est ou bien
l'identité, ou bien une « symétrie » --- mais alors elle échange les deux
$\frac12$-espaces déterminés par $H$, donc elle ne pourrait respecter $C$,
OK.
\note{à gauche de la fin de la démonstration, une petite figure : un
parallélogramme (un plan) traversé par un segment oblique marqué d'un
point, l'hyperplan $H$ et une droite qui le perce}

\emph{Corollaire} Soit $C_0$ un cube de dim $n$, $D_0$ un drapeau dudit,
\struck{Al} $C$ un cube de dim $n$. Alors $u \mapsto u(D_0)$
\[
  \mathrm{Isom}(C_0, C) \xrightarrow{\ \sim\ } \mathrm{Drap}(C)
\]

\page{52}
donc $\mathrm{Drap}(C)$ est un torseur \emph{à droite} sous $G_0 =
\mathrm{Aut}(C_0)$ (et un torseur à gauche sous $G = \mathrm{Aut}(C)$, qui
est le \uncertain{commutant} de $G_0$\ldots)

Le foncteur
\[
  \begin{array}{ccc}
    C & \longmapsto & \mathrm{Drap}(C) \\
    n\text{-Cubes} & \longrightarrow & G_0\text{-tors.\ à dr.}
  \end{array}
\]
est une équivalence de catégories.
\note{un double trait vertical dans la marge gauche accompagne la fin de la
page précédente et ces lignes ; la page commence au milieu de la phrase
du corollaire de la page 51}

\emph{NB} $(C_0, D_0)$ est défini à isom.\ unique près --- donc aussi
$G_0$. Si on veut, on prend pour $C_0$ le cube standard $C_n$,
\struck{qui \uncertain{dans} \ill{} \ill{}} \add{qui \ill{} \ill{}} de
une combinaison \uncertain{canonique} : $\mathbb{B} = [1, n] \times
\lbrace \pm 1 \rbrace$, avec comme drapeau celui défini par la section cte
de valeur 1, et ses restrictions aux sous-intervalles $[1, i]$ ($0
\leqslant i \leqslant n$). Ainsi on trouve un isom.\ can.
\[
  G_0 \simeq \lbrace \pm 1 \rbrace^n \cdot \mathfrak{S}_n \simeq
  \mathrm{Aut}(C_n)
\]
\uncertain{produit} ½ direct
\note{la ligne « qui … » est biffée d'un trait et récrite au-dessus, très
cursivement ; « ½ direct » : il écrit ½ pour « semi- »}

\page{53}
Symétrie des cubes : c'est l'homothétie \struck{\ill{}} $-\mathrm{id}$.
Du pt de vue combinatoire, c'est l'automorphisme $I' \mapsto \sigma I'$
dans l'ens.\ des « sections partielles » de $\mathbb{B}$ sur $I$. Du pt de
vue \uncertain{maquette} des cubes, c'est l'automorphisme $\sigma$ de
$(\mathbb{B}, \sigma)$. C'est l'unique automorphisme \add{\uncertain{un}}
\uncertain{laissant} le\ldots
\note{la phrase reste en suspens en fin de ligne}

\struck{\emph{Prop} Soient $F$, $F'$ deux facettes \add{fermées} non vides,
Alors $F \cap F' \neq \emptyset$ \emph{ou} $F \cap \sigma F' \neq
\emptyset$. (Ceci caractérise l'autom.\ $\sigma$ ?) \ill{} \ill{} Cl.\
revient à ceci : si $I'$, $I'' \subset \mathbb{B}$ sont deux sections
partielles de $\mathbb{B}$ sur $I = \mathbb{B}/\sigma$, alors \ill{}}
\note{tout ce bloc est encadré et barré de traits obliques ; « fermées »
est ajouté au-dessus de « non vides »}

\ldots\ \struck{\ill{}} propriété : si $F$ est une facette de codim 1
(une hyperfacette), alors $\sigma F$ est une hyperfacette
\uncertain{ne} rencontrant pas $F$. Plus généralement, pour une facette $F$
\struck{de dimension} \add{\ill{}} à part celle de dim $n$, $F$ et $\sigma
F$ ne se rencontrent pas (en fait les hyperplans d'appui ne se rencontrent
pas et \uncertain{sont} \ill{})
\note{le début de ce paragraphe est pris dans l'encadré biffé ; « \ill{}
propriété » paraît en être la reprise ; « de dimension » est biffé et
remplacé par un mot illisible}

\marginal{$/\!/$\,)}
Aussi : le centr.\ de $\mathrm{Aut}(C)$ est le groupe $\lbrace 1, \sigma
\rbrace$ [\emph{NB} si $n \geqslant 3$, le centre de $\mathfrak{S}_I$ est
$\lbrace 1 \rbrace$\ldots], donc $\sigma$ est l'unique élément central $\neq
1$.
\note{« $/\!/$ ) » : deux traits obliques et une parenthèse dans la marge,
devant la ligne qui finit par « ne se rencontrent pas » ; un trait
vertical accompagne le paragraphe « Aussi »}

\marginal{\ill{}}
Nb des facettes de dim $d$ ($0 \leqslant d \leqslant n$)
\[
  c_d = \binom{n}{d} 2^{n-d}
\]
On a la formule d'EP
\begin{align*}
  \sum_0^n (-1)^d c_d &= \sum (-1)^d \binom{n}{d} 2^{n-d} = (-1)^n \sum
  \binom{n}{d} 1^d (-2)^{n-d} \\
  &= (-1)^n (1-2)^n = 1 \qquad \text{i.e.}
\end{align*}
\[
  \boxed{\ \sum_0^n (-1)^d c_d = 1\ }
\]
(\emph{valable en fait pour tt polyèdre convexe de dim $n$} !)
\note{un mot court dans la marge gauche, devant une accolade qui embrasse
le calcul, illisible (le même que devant le corollaire de la page 50) ;
« EP » : Euler--Poincaré ; le signe $\leqslant$ devant le second $n$ de
« $0 \leqslant d \leqslant n$ » est repassé ; le $\sum$ de la seconde
somme est surchargé}

\page{54}
\struck{Ens.}\ d'orientations des cubes (ou de l'espace affine qu'il
engendre, ce sera provisoirement notre définition ici\ldots)
\[
  \omega_C \overset{\text{déf}}{=} \omega_E \simeq \omega_I \wedge
  \bigwedge_{i \in I} \mathbb{B}_i
\]
\note{au-dessus du premier signe, « déf » et un mot court,
\uncertain{« prov »} (provisoire) ; « Ens. » est surchargé en tête de
ligne}

L'effet de $\sigma$ sur $\omega_C$ est $(-1)^n$, donc

\emph{Prop} L'antipodisme conserve l'orientation ssi $n$ pair (donc la
renverse ssi $n$ impair)

\emph{Cor} Soit $n$ impair, $\mathrm{Aut}(C) \simeq \mathrm{Aut}^+(C)
\times \lbrace \pm 1 \rbrace$
\note{le $\simeq$ est repassé}

\emph{NB} (Thème de modélisation géométrique) : Soit $C$ un polyèdre
convexe \struck{\ill{}} \supplied{engendrant l'}espace affine $E$. Alors l'ens.\
à 2 él.\ $\omega_E$ peut se décrire canoniquement en termes des polyèdres
combinatoires $\mathrm{Comb}(C)$.

\emph{Indication} L'\struck{\ill{}} orientation de $E$ permet de définir
l'ens.\ \struck{$D(\omega)$} (des \add{repères} \struck{drapeaux} directs »
$F_0 \subset F_1 \subset \cdots \subset F_{n-1} \subset F_n = C$
(\add{\uncertain{des}} facettes) du polyèdre $C$, et elle est définie par
le sous-ens.\ de \add{$D(\omega)$} \struck{\ill{}}. Il faut donc
décrire, en termes de la seule structure d'ensemble ordonné de
$\Phi = \mathrm{Comb}(C)$, ce qu'on appellera les parties de
\add{$D(\Phi)$} \struck{$\mathrm{Comb}(C)$} (l'ens.\ des \add{repères}
\struck{drapeaux} de $\Phi$) qui sont des $D(\omega)$ [définition
combinatoire de l'orientation]. \emph{Solution} Soit $D^+ \subset D$,
$D^- = D - D^+$, on veut que \add{les} $\sigma_i$ ($1 \leqslant i
\leqslant n$) sur $D$ transforment $D^+$ en $D^-$ (donc $D^-$ en
$D^+$\ldots). [En fait, il y a exactement 2 telles partitions : $D^+$ et
$D^-$ !
\note{$\Phi$ est son I barré ; la parenthèse ouvrante devant « des
repères » se ferme avant « directs », le guillemet ouvrant manque ; le
« $D(\omega)$ » biffé à la première ligne de l'indication est ajouté
plus bas au-dessus d'un passage biffé illisible ; la dernière ligne est
serrée sous la précédente, et son crochet ne se ferme pas}

\page{55}
\emph{Drapeaux d'un polyèdre convexe --- drapeaux d'un ens.\ ordonné.
Repères d'un polyèdre convexe}
\note{une petite flèche vers la gauche au-dessus du titre}

\emph{Prop.} Soit $C$ un polyèdre convexe, $\mathrm{Fac}(C)$ le polyèdre
combinatoire associé, $\mathrm{Rep}(C) = \mathrm{Rep}(\mathrm{Fac}(C))$
l'ens.\ des repères. Alors les \struck{applications} homom.\ de groupes
\[
  \mathrm{Aut}(C) \xrightarrow{\ \alpha\ }
  \mathrm{Aut}_{\mathrm{ord}}(\mathrm{Fac}(C)) \xrightarrow{\ \beta\ }
  \mathrm{Aut}_{\mathrm{ens}}\, \mathrm{Rep}(C)
\]
sont injectifs. Plus précisément par $\beta$,
$\mathrm{Aut}_{\mathrm{ord}}(\mathrm{Fac}(C))$ opère \emph{librement} sur
$\mathrm{Rep}(C)$ (a fortiori, $\alpha$ étant injectif, $\mathrm{Aut}(C)$
opère librement)
\note{sous $\mathrm{Aut}(C)$, un « $\cdot H$ » peu clair ; le premier
indice de $\mathrm{Aut}$ au-dessus de $\mathrm{Rep}(C)$ est surchargé ;
« Fac » est lu comme à la page 49}

En effet, $\alpha$ est inj.\ car un auto.\ \uncertain{d'un} polyèdre
convexe $C$ est connu quand on connaît sa restriction aux pts extrémaux.

Soit $u$ un automorphisme de $\Phi = \mathrm{Fac}(C)$, fixant un drapeau
$(F_0, \dots, F_{n-1}, F_n)$. Par hyp.\ de récurrence, il \struck{fixe}
\struck{$P$\ill{}} est l'identité sur $\Phi_{\le F_{n-1}}$. D'autre
part, s'il est l'identité sur une \struck{facette} $\Phi_{\le x}$, $x$ une
facette de dim $n-1$, et si $y$ est une facette telle que $x \wedge y$ soit
de dim $n-2$, alors il est l'identité sur $\Phi_{\le y}$ --- car il fixe
un drapeau de $\Phi_{\le y}$. Il \add{\uncertain{faut voir}} \struck{\ill{}}
que le \uncertain{rel.\ d'équiv.}\ engendrée par « $x$ et $y$ ont une
facette de dim $n-2$ en commun » est la relation grossière. \struck{\ill{}
\ill{}} Argument topologique de connexité\ldots
\note{il écrit $\Phi_x$, $\Phi_y$ (son I barré, indice $x$ ou $y$), lu ici
$\Phi_{\le x}$ pour « la partie de $\Phi$ au-dessous de $x$ » ; au-dessus de
la ligne, « \struck{\ill{}}$_{n-1}$ » ; la fin de la phrase « Il
\ldots\ grossière » est très cursive}

\emph{Déf} $C$ \struck{\ill{}} est dit polyèdre \emph{régulier} si
$\mathrm{Aut}_{\mathrm{aff}}(C)$ est \uncertain{transitif} (\uncertain{donc}
simplement transitif) sur $\mathrm{Rep}(C)$, homogène si pour tt $0
\leqslant i \leqslant n-1$, $\exists\ d_i \in \mathbf{N}$ tel que $\forall F
\in \mathrm{Fac}_{i+1}(C)$, le cardinal des
\note{la page s'arrête sur « des » ; l'indice de $\mathrm{Aut}$ est lu
$\mathrm{aff}$ avec doute, et la borne $n-1$ est surchargée}

\page{56}
l'ens.\ des \struck{facettes} $y \in \mathrm{Fac}_{i}(C)$ qui majorent $x$
soit $d_i$. (Donner des ex.\ \ill{}.)
\[
  \mathrm{card}\, \mathrm{Rep}(C) = d_0 d_1 \cdots d_{n-1} = d_0 \cdots
  d_{n-1}
\]
(car $d_n = 1$). \struck{Dans ce cas, si $N = \mathrm{card}\,
\mathrm{Aut}(C)$, $N_c = \mathrm{card}\, \mathrm{Aut}(\mathrm{Fac}(C))$}
\note{un trait vertical dans la marge accompagne la fin de la définition
de la page 55 ; les indices $i+1$, $i$ sont peu lisibles ; dans le produit
central l'un des indices est surchargé}

Évidemment régulier $\Rightarrow$ homogène. \struck{Évidemment, \ill{}}
\struck{$N \mid N_c \mid \mathrm{card}\, \Phi$}
\note{« $N \mid N_c \mid$ card » est biffé et repris plus bas}

Soit
\begin{align*}
  N &= \mathrm{card}\, \mathrm{Aut}(C) \\
  N_c &= \mathrm{card}\, \mathrm{Aut}(\mathrm{Fac}(C)) \\
  R &= \mathrm{card}\, \mathrm{Rep}(C),
\end{align*}
on a
\[
  N \mid N_c \mid R
\]
et $C$ régulier $\Longleftrightarrow$ $N = R$ (donc $R = N_c$), dans les
autres cas \uncertain{combinatoires} \ill{} (\uncertain{même} \ill{}
\uncertain{dimension}). Dans ce cas on a \ldots

\uncertain{une} \uncertain{décomposition}
\[
  N = N_c = d_0 \cdots d_{n-1}
\]
\note{les deux lignes « dans les autres cas \ldots\ Dans ce cas on a
\ldots » sont très cursives ; un double trait horizontal sépare ce qui
précède de la suite}

\emph{Orientation d'un polyèdre} \add{convexe} $C$ : par définition, ce
sera une orientation de l'espace affine $E_C$ qu'il engendre.
\struck{\ill{} de ses \ill{} \ill{}}
(\struck{Mais on montre que} l'ens.\ $\omega(C)$ en dépend, à isom.\
canonique près, \uncertain{que} de $\Phi(C) = \mathrm{Fac}(C)$ (polyèdre
combinatoire strict associé)). Donc il y a \struck{\ill{}}
\[
  \mathrm{Isom}_{\mathrm{ord}}\bigl(\Phi(C), \Phi(C')\bigr) \longrightarrow
  \mathrm{Isom}\bigl(\omega(C), \omega(C')\bigr)
\]
\[
  \mathrm{Aut}_{\mathrm{ord}}\bigl(\Phi(C)\bigr) \longrightarrow \lbrace
  \pm 1 \rbrace = \mathrm{Aut}_{\mathrm{ens}}\bigl(\omega(C)\bigr) \qquad
  (\text{hom.\ de groupes})
\]
Cas où $C$ admet une antipodisme : sa sign.\ est $(-1)^n$
\note{« convexe » est ajouté au-dessus de la ligne ; le passage biffé
« Mais on montre que », au-dessus duquel un premier ajout est lui aussi
biffé, reste en partie illisible ; l'indice de $\mathrm{Aut}_{\mathrm{ens}}$
est surchargé}

\page{57}
\emph{Automorphismes des $n$-cubes.}

Considérons le $n$-cube de la maquette $(\mathbb{B}, \sigma)$, soit $I =
\mathbb{B}/\sigma$. Un automorphisme correspond : un automorphisme $u$ de
la maquette, \uncertain{induisant} un autom.\ $u_I$ de $I$, d'où une
décomposition de $I$ en orbites
\[
  I = \coprod_{\alpha} I_\alpha
\]
($u_I$ est connu à conjugaison près quand on connaît les cardinaux des
$I_\alpha$, rangés par ordre croissant\ldots) Soit $\mathbb{B}_\alpha$
l'image inverse de $I_\alpha$, donc $\mathbb{B}_\alpha$ est stable par $u$,
\struck{donc \ill{} d'orbites, il y a des \ill{}} \struck{\ill{} sur
$I_\alpha$ \ill{} \ill{} degré \ill{}. Donc} \struck{2 cas se
présentent :} \struck{1°) $\mathbb{B}_\alpha$ est une orbite. Soit
$n_\alpha$}
\note{ces lignes sont barrées ; à droite, reliée par une accolade :
« la situation est \uncertain{$\sim$} somme des
$(\mathbb{B}_\alpha, u_\alpha = u|\mathbb{B}_\alpha)$ sur $I_\alpha =
\mathbb{B}_\alpha / (\sigma|\mathbb{B}_\alpha)$ »}

et on est ramené à étudier chacun de ces cas.
\note{un mot court à la fin de la ligne, surchargé, illisible}
(Géométriquement, $C = \prod C_\alpha$, où $C_\alpha$ est le cube associé à
$(\mathbb{B}_\alpha, \sigma|\mathbb{B}_\alpha)$, et $u = \prod u_\alpha$,
\ldots). Cela nous ramène à étudier le cas où $u_I$ est transitif sur $I$.

$\mathbb{B}$ est réunion d'orbites, dont chacune s'envoie sur $I$
(l'unique orbite \struck{\ill{}} de $u_I$ dans $I$), \uncertain{avec} un
degré déterminé. Donc il y a 2 cas

a) Deux orbites isomorphes : $I \xrightarrow{\ \sim\ }$ par $p : \mathbb{B}
\to I$. Alors, si \struck{\ill{}} $\varepsilon$ est l'ens.\ des 2 orbites,
\ldots.\ on a canoniquement
\[
  \mathbb{B} \simeq \varepsilon \times I, \qquad u =
  \mathrm{id}_\varepsilon \times u_I
\]
\note{« a) » est lu sur un signe peu net ; dans « $I \xrightarrow{\sim}$
par $p$ », la flèche est un long trait sous lequel « par » est écrit ;
le $\mathrm{id}$ devant $u_I$ est biffé en surcharge ; la paire
$(\mathbb{B}_\alpha, \sigma|\mathbb{B}_\alpha)$ est lue sous une rature}

\page{58}
où $u_\varepsilon$ est l'\struck{\ill{}} permutation \ill{} identique de
$\varepsilon$ (toutes les fibres sont canoniquement isom.\ par les
itérés de $u$)

\marginal{« \uncertain{cas tordu} »}
b) Il y a une seule orbite dans $\mathbb{B}$. Cette situation est aussi
uniquement déterminée (\uncertain{connaissant} \add{$I \to$} \struck{\ill{}}
\uncertain{à isom.\ près} une fois $I$, $u_I$ donné, \uncertain{modulo}
\uncertain{isom.} une fois $n = \mathrm{card}\, I$ donné).
\[
  \begin{array}{c}
    \phantom{u b_n =}\ \overset{u b_1}{\|} \qquad\qquad
    \overset{u^{n-1} b_n}{\|} \\
    \begin{array}{r|cccc}
      & b_1 & b_2 & \cdots & b_n \\
      u b_n = & b'_1 & b'_2 & \cdots & b'_n \\
      \hline
      & i_1 & i_2 & \cdots & i_n
    \end{array} \\
    \phantom{u b_n = i_1}\ \underset{u i_1}{\|} \qquad\quad
    \underset{u^{n-1} i_1}{\|}
  \end{array}
\]
\note{tableau à deux lignes $b_1, \dots, b_n$ et $b'_1, \dots, b'_n$ (les
deux éléments de chaque fibre) au-dessus d'un trait, et $i_1, \dots, i_n$
au-dessous ; il note $b_2 = u b_1$, $b_n = u^{n-1} b_1$ (lu $u^{n-1}b_n$,
surchargé), $b'_1 = u b_n$, $i_2 = u i_1$, $i_n = u^{n-1} i_1$ ; la mise en
place des signes $=$ verticaux est approximative}

\marginal{\emph{NB} Classification des \uncertain{fibrés} $\mathbb{B} \to
\overline{I}$ ($\mathrm{card}\, I$) \uncertain{décrits} \uncertain{par}
\uncertain{$\mathbf{Z}$-ens.}\ à \uncertain{\ill{}} finis}
\note{note de la moitié droite de la page, séparée par un trait
vertical ; la barre sur $I$ est surchargée}

Ainsi, pour $n = \mathrm{card}\, I$ fixé, les classes d'isom.\ des
automorphismes d'un $n$-cube (i.e.\ les classes de conjugaison d'éléments de
\[
  \mathrm{Aut}(C_n) = \lbrace \pm 1 \rbrace^n \cdot \mathfrak{S}_n\,)
\]
sont données par 2 suites d'entiers \uncertain{naturels}
\[
  \left.
  \begin{array}{l}
    \alpha_1^+ \ \cdots\ \alpha_n^+ \in \mathbf{N} \\
    \alpha_1^- \ \cdots\ \alpha_n^- \in \mathbf{N}
  \end{array}
  \right| \quad
  \begin{array}{l}
    \alpha_i^+ \ (\alpha_i^-) = \text{nb des cycles de long.\ } i \\
    \text{de } u_I \text{ au-dessus desquels il y a} \\
    \text{un cycle de long.\ } i \ (\text{resp.\ } 2i)
  \end{array}
\]
satisfaisant
\[
  \sum_i i \underbrace{(\alpha_i^+ + \alpha_i^-)}_{\alpha_i} = n
\]
\note{le $\alpha_i^+$ de la colonne de droite est surchargé ;
« (resp.\ $2i$) » est lu \uncertain{« (resp.\ 2i) »} ; le symbole devant
$i(\alpha_i^+ + \alpha_i^-)$ est un $\Sigma$ surmonté d'une barre}

\emph{Étude des facettes fixes} ($\neq \emptyset$) \struck{\ill{}}
Correspondent aux parties non vides \add{$\mathbb{B}'$} de $\mathbb{B}$ qui
a) stables par $I$ b) sections partielles. Au-dessus de chaque $I_\alpha$
correspondant au \uncertain{cas} \uncertain{tordu},
\note{la page s'arrête sur une virgule ; « a) stables par $I$ » tel
qu'écrit (on attend « par $u$ »)}

\page{59}
$\mathbb{B}'$ doit être vide, au-dessus de chaque $I_\alpha$ non tordu, il
est vide ou une des deux orbites de $\mathbb{B}_\alpha$ sur $I_\alpha$
(trois cas), de sorte que le nb des \add{\uncertain{telles}} facettes
\struck{\ill{}} (y compris la facette $C$) \struck{\ill{}} est
\[
  3^{\sum \alpha_i^+}
\]
\note{devant $3^{\sum \alpha_i^+}$, une formule biffée de plusieurs
traits, illisible}

Les facettes de codimension donnée $d$ correspondent aux parties
$\mathbb{B}'$ telles que \struck{\ill{}} $\mathrm{card}\, \mathbb{B}' = d$.
On trouve, si $f_d$ est le nb de ces facettes
\[
  \boxed{\ \sum_{i=0}^n f_d\, t^d = \prod_{i=1}^n (1 + 2t^i)^{\alpha_i^+}\ }
\]
\note{la somme porte l'indice $i = 0$ tel qu'écrit, pour $d$}

\marginal{$+, 8$ \quad $-, 8$}
\emph{Exemple} Automorphismes de $C_3$
« indécomposables »
\[
  \left\lbrace
  \begin{array}{lll}
    1^\circ)\ \text{non tordu :} & e_1 \to e_2 \to e_3 \to e_1 & (+) \\
    2^\circ)\ \text{tordu} & e_1 \to e_2 \to e_3 \to -e_1 & (-)
  \end{array}
  \right.
\]
\note{les signes $(+)$ et $(-)$ sont cerclés, en marge droite ; dans la
marge gauche, devant un trait vertical, « $+, 8$ » et « $-, 8$ » l'un
sous l'autre}

Le premier est une rotation \add{d'ordre 3} autour de la diagonale
joignant $e_1 + e_2 + e_3$ et $-e_1 - e_2 - e_3$. Il y a (sans compter la
facette pleine) 2 \emph{facettes} fixes, savoir les deux sommets
précédents.

Le deuxième est d'ordre 6, c'est \add{\uncertain{(mod.\ conj.)}}
\struck{\ill{}} le produit du précédent avec l'antipodisme [\emph{NB}
cette relation entre cas tordu et non tordu vaut chaque fois que la dim $n$
est impair]
\note{« (mod.\ conj.) » est écrit au-dessus de la ligne et relié par un
trait à « le » ; « l'antipodisme » est écrit sur un premier mot surchargé}

b) Autom.\ décomposables : deux cycles (dans $I$)
\struck{$\alpha_1^+\ \alpha_1^+\ \alpha_2$\ill{}}
\note{la page s'achève sur une formule biffée de hachures, illisible}

\page{60}
\marginal{$-, 6$ \quad $+, 6$ \quad $+, 6$ \quad $-, 6$}
\begin{itemize}
  \item[1°)] Non tordus tous deux \qquad $e_1 \to e_1$, $e_2 \to e_3$, $e_3
  \to e_2$ ---
  symétrie autour plan diagonal (\emph{ordre 2})
  \item[2°)] 1 tordu, 2 non tordu \qquad $e_1 \to -e_1$, $e_2 \to e_3$,
  $e_3 \to e_2$ ---
  symétrie autour, \uncertain{composée} du préc.\ avec antip.\
  (\emph{ordre 2})
  \item[3°)] 1 non tordu, 2 tordu \qquad $e_1 \to e_1$, $e_2 \to e_3$,
  $e_3 \to -e_2$ \quad (ordre 4) ---
  rotation d'un quart de tour autour de l'axe perpendiculaire : \ill{}
  face (\uncertain{axe}-face)
  \item[4°)] 1 tordu, 2 tordu \qquad $e_1 \to -e_1$, $e_2 \to e_3$, $e_3
  \to -e_2$ \quad (ordre 4) ---
  (conjugué au précédent, suivi de \struck{\ill{}} \add{antipodisme})
\end{itemize}
\note{dans la marge gauche, devant un trait vertical, un signe et un
nombre en face de chaque cas : « $-, 6$ », « $+, 6$ », « $+, 6$ »,
« $-, 6$ », comme à la page 59 (« $+, 8$ », « $-, 8$ ») et plus bas ;
la ligne « symétrie autour, composée du préc.\ avec antip. » est écrite
sur une première ligne biffée et récrite, en partie illisible}

c) Autom.\ décomposables à trois cycles (déterminé par $\alpha_1^+$ et
$\alpha_1^-$, $\alpha_1^+ + \alpha_1^- = 3$)

\marginal{$+, 1$ \quad $-, 3$ \quad $+, 3$ \quad $-, 1$}
\begin{itemize}
  \item[1°)] $\alpha_1^+ = 3$ \qquad identité
  \item[2°)] $\alpha_1^+ = 2$ \qquad $e_1 \to e_1$, $e_2 \to e_2$, $e_3 \to
  -e_3$ \quad symétrie par rapport à un plan \struck{\ill{}} coord.
  \item[3°)] $\alpha_1^+ = 1$ \qquad $e_1 \to e_1$, $e_2 \to -e_2$, $e_3
  \to -e_3$ \quad sym.\ par rapport à un axe coord.
  \item[4°)] $\alpha_1^+ = 0$ \qquad antipodisme
\end{itemize}
\note{les deux « symétries » des cas 2°) et 3°) sont écrites à droite,
reliées aux formules par un trait courbe ; « par rapport à » est une
lecture \uncertain{incertaine}}

En tout il y a 10 classes ? \quad ($10 = 2 + 4 + 4$)
\note{le « 0 » de « 10 » est repassé sur un premier chiffre}

\note{un trait horizontal barre la page ; ce qui suit est la « Prop »
dont la page 61 tire deux corollaires}

\emph{Prop} Soit $u$ un automorphisme de $C$ tel que pour tout sommet $x$,
on ait $ux = \pm x$. Alors \add{(si $n \geqslant 3$)} $u$ est l'identité
ou l'antipodisme.

Donc
\[
  \mathrm{Aut}(C)/(\pm 1) \hookrightarrow
  \mathrm{Aut}\bigl(\Phi_0(C)/\sigma\bigr)
\]
donc si $C$ de dim impaire
\[
  \mathrm{Aut}^+(C) \hookrightarrow
  \mathrm{Aut}\bigl(\Phi_0(C)\bigr)
\]
\note{un trait vertical dans la marge accompagne l'énoncé ; « (si $n
\geqslant 3$) » est écrit au-dessus de la ligne et relié à « $u$ est » ;
$\Phi$ est son I barré ; sous la première flèche, un court mot suivi d'un
signe $=$ doublement souligné, « l'$=$ », peu clair}

\end{document}
